\section*{\centering Abstract}
\phantomsection
\addcontentsline{toc}{section}{Abstract}

\vspace{1cm}
We investigate the morphological classification of extended radio sources using self-supervised learning. The analysis covers \(32\,514\) LoTSS-DR2 sources with largest angular sizes of \(60''\)--\(1000''\). Contrastive pre-training, manual cluster curation and a \(50\)-member classifier ensemble produce an eleven-class catalogue with morphology assignments and uncertainty estimates.

We introduce the similarity-to-radio-galaxy score \(S_\mathrm{RG}\), which combines support for radio-galaxy-like classes with ensemble reliability. The score distinguishes broad morphological resemblance from concentrated and consistent predictions, providing an additional criterion for selecting sources for visual inspection.

Comparisons with three external catalogues reveal stronger FRII than FRI correspondence in two catalogues, with similar performance in the third. Reference FRI sources frequently receive Bright-core or Structured-Bent assignments, highlighting the importance of central brightness and curvature in the learned classification. Application to ten manually isolated radio relics produces varied morphology assignments and low or transitional scores, illustrating how the framework describes diffuse emission through structures learned from radio galaxies.
